% ============================================================================
% Foreword --- Technical Reference
% ============================================================================

\section*{Foreword}
\addcontentsline{toc}{section}{Foreword}

This Technical Reference is the engineering counterpart of the
companion report \emph{Infrastructure and Organisation of
Information Systems for Universities --- Analysis, State of the Art
and Propositions}. The companion report makes the institutional
argument: it diagnoses the present state of university IT, surveys
the state of the art, and proposes a direction. The present
document does the architectural and operational work that
direction implies, at the level of detail a practitioner needs in
order to act.

\paragraph{Place of this document in the suite.} The suite of
which this Reference is part includes the policy-and-analysis
main report with its synthesis-oriented executive companion for
governance, an argumentation module addressing anticipated
objections, a regulatory-applicability annex, and the operational
summary that condenses the present Reference. The Technical
Reference is the artefact that travels independently to the
international university network anticipated in phase~2: it
contains the contracts, instantiations, alternatives, waves and
conformance material on which a peer institution can build its
own implementation without re-deriving the structural choices.

\paragraph{Posture.} The Reference adopts a deliberately stronger
voice than the rest of the dossier. It is normative about
\emph{interface contracts} (what every layer must expose for the
architecture to compose), illustrative about \emph{reference
instantiations} (the open-foundation defaults shown to be
sufficient), analytical about \emph{alternatives} (commercial or
hybrid), and conditional about institutional projections (where
the state of the art does not yet support a definite claim). It
is vendor-neutral by construction: open-source projects appear as
the demonstrable default because their source visibility and
standard interfaces make the contracts auditable, and proprietary
components are admitted as documented exceptions evaluated through
the sovereignty grid of \trchap{ch:sovereignty-grid}{}.

\paragraph{Authorship and contact.} The Reference is published
under the auspices of \keyword{SUIT} --- the
\emph{International Working Group on Sustainable University IT in
the AI Era} --- with the adopting institution's designated lead
author as initial contact point. The
group is intended to grow as peer institutions adopt and adapt
the architecture; contributions, corrections and counter-examples
from member institutions are expected to feed successive
revisions.

\paragraph{Always-living document.} This Reference is explicitly
an always-living document. It records the working state of an
ongoing engineering and governance investigation rather than a
finished doctrine. Successive versions will integrate the
operational experience of adopting institutions, the corrections
that peer review surfaces, and the evolution of the standards and
projects on which the architecture rests. Readers are encouraged
to engage critically with the material and to communicate their
findings: a supportive and cooperative posture by the readership
is the mechanism by which the document's quality is expected to
improve over time.

\paragraph{On AI assistance.} The engineering of this document
was partially supported by digital tools, including some that
incorporate artificial-intelligence features, used as drafting
aids and as cross-checking instruments. As stated on the title
page, the document may consequently contain imprecisions,
omissions, errors, or hallucinated material. A post-redaction
verification checklist is maintained in the working repository
and is applied at the close of each chapter; the institutional
precedents and external references catalogue
(\trchap{ch:institutional-precedents}{} and
\trchap{ch:external-references}{}) carry an explicit invitation
for member institutions to confirm and refine the claims made
about them.
